% !PRINTED 281
\weberpage{281}

%<<CONT
woraus man durch Zusammensetzung

\weq{16}{\psi' \Theta' \chi' = \begin{pmatrix} 0, & 1 \\ -1, & 0 \end{pmatrix}}

erhält. Wir können demnach in den drei Fällen folgende erzeugende Substitutionen der Tetra\"eder-, Octa\"eder- und Ikosa\"edergruppe [(6), (13), (15)] annehmen:

\[
p = 5. \quad \begin{pmatrix} 2, & 0 \\ 0, & 3 \end{pmatrix}, \ \begin{pmatrix} 0, & 1 \\ -1, & 0 \end{pmatrix}, \ \begin{pmatrix} 2, & 2 \\ 1, & -1 \end{pmatrix}
\]
\[
p = 7. \quad \begin{pmatrix} 3, & 0 \\ 0, & 5 \end{pmatrix}, \ \begin{pmatrix} 0, & 1 \\ -1, & 0 \end{pmatrix}, \ \begin{pmatrix} 1, & -2 \\ -1, & 3 \end{pmatrix}
\]
\[
p = 11. \quad \begin{pmatrix} 3, & 0 \\ 0, & 4 \end{pmatrix}, \ \begin{pmatrix} 0, & 1 \\ -1, & 0 \end{pmatrix}, \ \begin{pmatrix} -3, & 1 \\ 1, & 3 \end{pmatrix}.
\]

Durch wiederholte Zusammensetzung, die sich auf sehr verschiedene Arten anordnen lässt, ergeben sich hieraus leicht die vollständigen Gruppen:

\[
\begin{gathered}
p = 5.\\[2pt]
\begin{pmatrix} x, & 0 \\ 0, & x^{-1} \end{pmatrix}, \ \begin{pmatrix} 0, & x \\ -x^{-1}, & 0 \end{pmatrix}, \ \begin{pmatrix} x, & y \\ 2y^{-1}, & 3x^{-1} \end{pmatrix} \quad \begin{aligned} x &= 1, 2\\ y &= \pm 1, \pm 2 \end{aligned}
\end{gathered}
\]
\[
\begin{gathered}
p = 7.\\[2pt]
\begin{pmatrix} x, & 0 \\ 0, & x^{-1} \end{pmatrix}, \ \begin{pmatrix} 0, & x \\ -x^{-1}, & 0 \end{pmatrix}, \ \begin{pmatrix} x, & -xz \\ -2x^{-1}z^{-1}, & 3x^{-1} \end{pmatrix}, \ \begin{pmatrix} xz, & x \\ -3x^{-1}, & -2x^{-1}z^{-1} \end{pmatrix}\\[2pt]
x = 1, 2, 3; \quad z = 1, 2, 4\\
(z \text{ quadratischer Rest von } 7).
\end{gathered}
\]
\[
\begin{gathered}
p = 11.\\[2pt]
\begin{pmatrix} x, & 0 \\ 0, & x^{-1} \end{pmatrix}, \ \begin{pmatrix} 0, & x \\ -x^{-1}, & 0 \end{pmatrix}, \ \begin{pmatrix} x, & -xz \\ -x^{-1}z^{-1}, & 2x^{-1} \end{pmatrix}, \ \begin{pmatrix} xz, & x \\ -2x^{-1}, & -x^{-1}z^{-1} \end{pmatrix}\\[2pt]
x = 1, 2, 3, 4, 5; \quad z = 1, 3, 4, 5, 9\\
(z \text{ quadratischer Rest von } 11).
\end{gathered}
\]

Die Anwendung der Transformation (14) führt bei $p = 5$ zu keiner neuen Gruppe; bei $p = 7, 11$ erhält man je eine andere Gruppe, die aus diesen hervorgeht, wenn man für $z$ die Reihe der quadratischen Nichtreste statt der Reste setzt\textsuperscript{1}).

\bigskip
\hrule
\bigskip

\textsuperscript{1}) Die Eigenschaft der Congruenzgruppen, einfach zu sein, hat schon \gesperrt{Galois} gekannt. Ebenso waren ihm die Theiler vom Index $p$ für $p = 5, 7, 11$ bekannt. Eingehender untersucht sind diese Gruppen von
%CONT>>
